\section{项目任务与技术路线}
\label{sec:task}

输入「霜、孤舟」时，系统要在短小的篇幅内安置两个意象，写成四句五字。课程要求覆盖从模型搭建、预训练、微调到部署的完整过程。选择古诗是因为字数、句数、韵脚和字面用词可以复核，切题、可读性与意境又需要读原诗。这样，同一次实验可以检验程序是否达成任务约定，也能观察自动指标以外的生成效果。

Orchid 实现三个规模的字符级 GPT。14M 用于跑通基线和检查训练程序，42M 与 92M 用于同预算对照。最终系统以 92M 唐宋模型为起点，经监督微调和同韵诗续训，再在推理时按需限制字数、偶句韵部与用词。系统提供自由模式，使训练所得能力与推理时额外保障的规则可以在同一权重上分开测量。图~\ref{fig:workflow} 给出完整流程和验证所在的位置。

\reportfigure{workflow.pdf}{项目流程。公开语料先分组再派生分词窗口及监督样本。冻结唐诗验证集用于选模，独立新题用于报告生成表现。}{fig:workflow}

\begin{table}[!htbp]
  \centering
  \caption{报告覆盖 \texttt{Problem.md} 的项目评分要求}
  \label{tab:rubric}
  \begin{tabularx}{\linewidth}{@{}p{.23\linewidth}p{.10\linewidth}Y@{}}
    \toprule
    评分项 & 占比 & 对应内容与证据\\
    \midrule
    知识与原理 & 35\% & 自回归、因果注意力、位置嵌入、参数量推导、分词与监督目标，见第~\ref{sec:model}、\ref{sec:corpus} 和 \ref{sec:sft} 节。\\
    模型实现 & 20\% & 可训练的解码器、三种配置、同预算实验、单卡容量检查与部署，见第~\ref{sec:model}、\ref{sec:pretrain} 和 \ref{sec:deployment} 节。\\
    预训练 & 20\% & 固定快照、训练窗、随机初始化、优化日志与续训、学习曲线和原诗续写，见第~\ref{sec:corpus} 和 \ref{sec:pretrain} 节。\\
    微调 & 20\% & 监督标签、掩码、选模、韵律对照与生成测试，见第~\ref{sec:sft} 至 \ref{sec:evaluation} 节。\\
    成果与体会 & 5\% & 交付范围、失败样例、限制、复现记录与个人项目体会，见第~\ref{sec:deployment} 和 \ref{sec:conclusion} 节。\\
    \bottomrule
  \end{tabularx}
\end{table}

\subsection{已有工作与本项目取舍}

Transformer 以注意力为核心，GPT-2 展示了单向语言模型的自回归学习路线~\cite{vaswani2017attention,radford2019language}。中文古诗自动生成早期已有循环网络方案~\cite{zhang-lapata-2014-chinese}。九歌将押韵和诗体检查接入生成系统~\cite{guo-etal-2019-jiuge}，CharPoet 研究逐字生成对格式控制的作用~\cite{yu-etal-2024-charpoet}。本项目使用随机初始化的字符模型和固定唐宋诗料。论文结果用来解释设计选择，正文中的各组数值均来自 Orchid 自己的协议和实测。

\section{模型结构与学习目标}
\label{sec:model}

\subsection{自回归概率与目标函数}

设诗文经分词后的序列为 $x_0,\ldots,x_T$。训练时将输入与标签错开一个位置。第 $t$ 个位置可看见 $x_{\le t}$，输出对 $x_{t+1}$ 的未归一化分数。令 $\mathcal I$ 为参与计损失的位置集合，预训练的字词预测与监督微调可统一写成
\begin{equation}
    \mathcal L =
    -\frac{1}{|\mathcal I|}\sum_{t\in\mathcal I}
    \log p_\theta(x_{t+1}\mid x_{\le t}),
    \qquad \mathrm{PPL}=\exp(\mathcal L).
    \label{eq:nll}
\end{equation}
代码用交叉熵直接接收 logits，将批次的损失和按有效目标数汇总。有效目标数对不同长度的样本和最后一个小批次都能正确计权。训练借助已知前缀并行预测各位置，运行时则每次从最后位置选一个标记并接回序列，遇结束符或达到预算后停止。监督阶段和预训练阶段的目标集合不同，比较 NLL 时始终在同一口径和相同留出集内进行。

\subsection{嵌入、因果注意力与前馈块}

每个位置将字符嵌入与可学习的绝对位置嵌入相加。对输入 $X\in\mathbb R^{T\times d}$，每个注意力头分别由线性投影产生 $Q,K,V$。预归一化块以 LayerNorm 后的表示计算
\begin{equation}
  \operatorname{Attn}(Q,K,V)
   = \operatorname{softmax}\left(\frac{QK^{\mathsf T}}{\sqrt{d_h}}+M\right)V,
  \qquad
  M_{ij}=
    \begin{cases}
      0,&j\le i,\\
      -\infty,&j>i.
    \end{cases}
\end{equation}
其中 $d_h=d/H$，用 $\sqrt{d_h}$ 缩放点积，可以让头宽变化时的注意力分数保持在较适合 softmax 的量级。掩码让位置 $i$ 只从自己和此前位置取值。将各头结果拼接并投影后，与块输入做一次残差加法。第二个子层使用 $d\to4d\to d$ 的 GELU 前馈网络，并进行另一次残差加法。末端归一化后，用词嵌入矩阵的转置生成词表 logits。

\reportfigure{transformer.pdf}{92M 模型结构和因果掩码。左侧的注意力与前馈组成一个块，重复 12 次。右侧绿色单元允许当前位置访问对应键。}{fig:transformer}

实现提供 PyTorch 的 scaled dot-product attention 与显式掩码路径，并在评估模式下关闭 dropout。测试检查两条注意力路径的输出差异、未来 token 对当前 logits 的隔离，以及权重保存、加载和训练续跑。共用输入嵌入与输出头的权重后，含偏置的参数量为
\begin{equation}
    N=d(V+C)+L(12d^2+13d)+2d,
\end{equation}
其中 $V$ 是词表大小，$C$ 是最大上下文，$L$ 是块数。块内的注意力和前馈分别贡献 $4d^2$ 与 $8d^2$ 个矩阵参数，两组 LayerNorm 及各线性层偏置贡献 $13d$ 个参数。模型共享 $V=8{,}582$、$C=256$、头宽 64 与 0.1 的训练 dropout。代码在实例化后逐张量计数核对配置。

\begin{table}[!htbp]
\centering
\caption{Orchid 的三种模型配置。实际参数数包含嵌入与归一化参数。}
\label{tab:model}
\begin{tabular}{@{}lrrrrr@{}}
\toprule
版本 & 块数 & 隐藏宽度 & 注意力头 & 前馈宽度 & 实际参数\\
\midrule
14M & 6 & 384 & 6 & 1,536 & 14,041,344\\
42M & 12 & 512 & 8 & 2,048 & 42,354,688\\
92M & 12 & 768 & 12 & 3,072 & 91,843,584\\
\bottomrule
\end{tabular}
\end{table}

\section{语料、冻结划分与字符词表}
\label{sec:corpus}

\subsection{唐诗基线与宋诗扩展}

唐宋正文来自 \texttt{chinese-poetry} 仓库的同一固定提交，唐诗 58 个分片，宋诗 255 个分片~\cite{chinesepoetry}。预处理从 \code{paragraphs} 提取正文，做 Unicode NFC 归一化，凭句读恢复诗句，将标点规范为逗号与句号交替。原始文字和来源索引保留供复查。无句子、缺字标记和混入非汉字内容的记录被登记。标题及作者用于追溯来源。监督目标必须本身可解析为四句五字，长诗只用于预训练与重合检查。

唐诗 57,607 条原始记录有 51,945 条通过初筛，按正文去重后为 51,138 首。按诗作分组，在相同或可检测的近重复关系上建立连边，再做约 90/5/5 的固定划分。训练、验证、测试分别有 46,031、2,557、2,550 首。近重复规则检查长诗内嵌的四句、四句诗三句对齐及有限编辑距离，让相关文本尽量同组。训练和评估窗口继承原有划分。这个边界特别重要，逐行随机分配会把同一首诗的前后句分到不同集合。

\reportfigure{corpus.pdf}{语料清洗与训练集扩容。右侧新增的宋诗只进训练集，唐诗验证与测试的诗文、标记与监督样本均按指纹核验。}{fig:corpus}

扩容时对宋诗重复执行正文清洗，筛除唐诗同文和与留出诗符合近重复规则的候选。宋诗中清洗、去重及泄漏筛查后保留 250,790 首，新训练集合达 296,821 首，词流含 18,451,787 个标记。宋诗加入后仍使用唐诗训练集学得的字表，以维持模型结构与留出集口径；扩展训练正文中的未登录 token 约占 0.374\%。该规则可以追踪已知重合，异文和典故借用仍可能穿过规则边界。

\subsection{字符级分词与两类训练样本}

字符词表仅由唐诗训练集拟合，固定顺序为保留标记加按字符排序的正文与标点。保留标记包括序列开头、结尾、填充、未知字，以及专用诗体和关键词分隔符。以字符为单位便于按汉字数施加规则，第 2、4 句的位置也可以落实到确定的生成步。代价是古诗很短时仍需逐字解码，生僻字和简繁异体须受词表覆盖约束。界面保留训练采用的字形，并对未登录字符返回明确提示。

预训练将分好的诗文流按上下文窗切片，每个位置预测其后标记。完整窗有 256 个位置，末尾填充位置的标签屏蔽。监督阶段从确认为五言四句的诗作中构造「关键词 $\to$ 诗文」对。唐诗与唐宋版训练监督集分别为 2,929、12,394 对，两版共有同一组 166 条验证和 157 条测试配对。与写作者给出的题意相比，训练标签只是可定位的弱监督。独立诗题、语义评分与显式约束分别检查条件创作的三个侧面。
